\documentclass[11pt,a4paper]{article}

% ---------------------------------------------------------------------------
% Ganti dengan identitas kamu
% ---------------------------------------------------------------------------
\newcommand{\Penulis}{Nama Penulis}
\newcommand{\Institusi}{Program Studi / Universitas}

\usepackage[utf8]{inputenc}
\usepackage[T1]{fontenc}
\usepackage[indonesian]{babel}
\usepackage{lmodern}
\usepackage{microtype}
\usepackage[a4paper,margin=2.5cm]{geometry}
\usepackage{amsmath,amssymb}
\usepackage{graphicx}
\usepackage{booktabs}
\usepackage{tabularx}
\usepackage{float}
\usepackage[font=small,labelfont=bf]{caption}
\usepackage{enumitem}
\usepackage{xcolor}
\usepackage{listings}
\usepackage{fancyhdr}
\usepackage[round,authoryear]{natbib}
\usepackage[hidelinks]{hyperref}

\graphicspath{{generated/}}
\input{generated/angka.tex}

\definecolor{biru}{HTML}{2A78D6}
\definecolor{abu}{HTML}{52514E}
\definecolor{latarkode}{HTML}{F6F6F4}
\definecolor{peringatan}{HTML}{C93A39}

\lstset{
  language=Python,
  basicstyle=\ttfamily\scriptsize,
  keywordstyle=\color{biru}\bfseries,
  commentstyle=\color{abu}\itshape,
  stringstyle=\color{teal!70!black},
  backgroundcolor=\color{latarkode},
  numbers=left,
  numberstyle=\tiny\color{abu},
  numbersep=6pt,
  frame=single,
  rulecolor=\color{black!15},
  breaklines=true,
  breakatwhitespace=true,
  showstringspaces=false,
  columns=fullflexible,
  keepspaces=true,
  tabsize=4,
  xleftmargin=1.5em,
  framexleftmargin=1.5em,
}

\setlength{\parskip}{0.5em}
\setlength{\parindent}{0pt}
\setlist{itemsep=0.2em}

\pagestyle{fancy}
\fancyhf{}
\fancyhead[L]{\small Prediksi Return Harian ASII, BBRI, dan TLKM 2026}
\fancyhead[R]{\small\thepage}
\renewcommand{\headrulewidth}{0.4pt}

\newcommand{\peringatansintetis}{%
  \ifsintetis
  \begin{center}
  \fcolorbox{peringatan}{peringatan!6}{\parbox{0.92\linewidth}{\small
    \textbf{\color{peringatan}PERHATIAN: DATA SINTETIS.} Seluruh angka, tabel, dan grafik
    dalam versi laporan ini dibuat dari harga acak (\emph{random walk}) untuk menguji
    program, \emph{bukan} data pasar sebenarnya. Jalankan
    \texttt{python -m stockpred.report --excel data/data\_saham.xlsx} untuk membuat versi
    dengan data asli.}}
  \end{center}
  \fi}

\begin{document}

% ---------------------------------------------------------------------------
% Halaman judul
% ---------------------------------------------------------------------------
\begin{titlepage}
  \centering
  \vspace*{3cm}
  {\Large\bfseries Prediksi Return Harian Saham\\[0.3em]
   Astra International (ASII), Bank Rakyat Indonesia (BBRI),\\[0.3em]
   dan Telkom Indonesia (TLKM) Selama Tahun 2026\par}
  \vspace{0.8cm}
  {\large Perbandingan Model Baseline, Regresi Ridge, dan LightGBM\\
   dengan Validasi \emph{Walk-Forward}\par}
  \vspace{2cm}
  {\large \Penulis\par}
  \vspace{0.3cm}
  {\Institusi\par}
  \vspace{0.3cm}
  {\TanggalDibuat\par}
  \vfill
  \peringatansintetis
  \vspace{1cm}
\end{titlepage}

% ---------------------------------------------------------------------------
\begin{abstract}
\noindent
Laporan ini menguji apakah return harian tiga saham berkapitalisasi besar di Bursa Efek
Indonesia, yaitu ASII, BBRI, dan TLKM, dapat diprediksi selama tahun 2026. Data harga harian
sejak \DataAwal{} diolah menjadi fitur teknikal seperti return masa lalu, volatilitas, RSI,
MACD, dan volume. Dua model \emph{machine learning}, yaitu regresi Ridge dan LightGBM,
dibandingkan dengan tiga model baseline sederhana. Evaluasi dilakukan secara
\emph{out-of-sample} dengan validasi \emph{walk-forward}: model hanya dilatih dengan data
sebelum hari yang diprediksi dan dilatih ulang setiap bulan. Prediksi juga diuji melalui
simulasi strategi perdagangan dengan biaya transaksi. \RingkasanHasil{}
\end{abstract}

\tableofcontents
\newpage

\peringatansintetis

% ---------------------------------------------------------------------------
\section{Pendahuluan}

\subsection{Latar Belakang}
Prediksi harga saham adalah salah satu topik paling populer dalam penerapan \emph{machine
learning}. Banyak proyek melaporkan akurasi yang sangat tinggi, tetapi sering kali hasil
tersebut berasal dari kesalahan metodologi. Kesalahan pertama adalah memprediksi
\emph{level} harga, bukan perubahannya. Karena harga besok hampir selalu dekat dengan harga
hari ini, model yang hanya menyalin harga kemarin pun akan terlihat sangat akurat di grafik.
Kesalahan kedua adalah kebocoran data (\emph{data leakage}), misalnya membagi data latih dan
data uji secara acak sehingga model ``melihat'' masa depan. Kesalahan ketiga adalah tidak
membandingkan model dengan baseline sederhana dan tidak memperhitungkan biaya transaksi.

Menurut hipotesis pasar efisien \citep{fama1970}, harga saham sudah mencerminkan informasi
yang tersedia, sehingga perubahan harga harian sulit diprediksi. Laporan ini menguji klaim
tersebut secara empiris pada tiga saham Indonesia dengan metodologi yang menghindari
kesalahan-kesalahan di atas.

\subsection{Rumusan Masalah}
\begin{enumerate}
  \item Apakah model \emph{machine learning} dapat memprediksi arah dan besar return harian
        ASII, BBRI, dan TLKM selama tahun 2026 lebih baik daripada model baseline sederhana?
  \item Apakah prediksi tersebut menghasilkan prediksi harga yang lebih akurat daripada
        prediksi naif ``harga besok sama dengan harga hari ini''?
  \item Apakah strategi perdagangan berdasarkan prediksi model menghasilkan kinerja lebih
        baik daripada membeli dan menahan saham (\emph{buy \& hold}) setelah biaya transaksi?
\end{enumerate}

\subsection{Tujuan}
Membangun dan mengevaluasi model prediksi return harian dengan validasi yang jujur, lalu
menilai manfaat praktisnya melalui simulasi strategi perdagangan.

\subsection{Batasan}
\begin{itemize}
  \item Saham yang diteliti: ASII, BBRI, dan TLKM.
  \item Periode uji: seluruh hari bursa tahun 2026 sampai \UjiAkhir.
  \item Horizon prediksi: satu hari bursa ke depan, menggunakan harga penutupan.
  \item Fitur hanya berasal dari data harga dan volume, tanpa data fundamental atau berita.
\end{itemize}

% ---------------------------------------------------------------------------
\section{Data}

Data harga harian diunduh dari Yahoo Finance menggunakan paket Python \texttt{yfinance}
(kode \texttt{ASII.JK}, \texttt{BBRI.JK}, dan \texttt{TLKM.JK}) dan disimpan ke file Excel
dengan satu sheet per saham (Lampiran~\ref{sec:kode}). Setiap baris berisi harga
pembukaan (\emph{Open}), tertinggi (\emph{High}), terendah (\emph{Low}), penutupan
(\emph{Close}), penutupan yang disesuaikan (\emph{Adj Close}), dan volume.

Harga yang dipakai dalam analisis adalah harga yang disesuaikan dengan dividen dan
\emph{stock split}. Tanpa penyesuaian, harga akan ``turun'' pada tanggal ex-dividen padahal
investor tidak rugi, sehingga return yang dihitung menjadi keliru. Baris dengan volume nol
(hari libur atau suspensi) dibuang.

\begin{table}[H]
  \centering
  \caption{Statistik deskriptif data. Return dan volatilitas dihitung dari log return harian.}
  \label{tab:statistik}
  \small
  \resizebox{\linewidth}{!}{\input{generated/tabel_statistik.tex}}
\end{table}

\begin{figure}[H]
  \centering
  \includegraphics[width=\linewidth]{harga.pdf}
  \caption{Harga penutupan yang disesuaikan. Area berwarna adalah periode uji (tahun 2026).}
  \label{fig:harga}
\end{figure}

% ---------------------------------------------------------------------------
\section{Metodologi}

\subsection{Variabel Target}
Yang diprediksi adalah \emph{log return} satu hari ke depan, bukan harga:
\begin{equation}
  r_{t+1} = \ln\!\left(\frac{P_{t+1}}{P_t}\right),
\end{equation}
dengan $P_t$ harga penutupan hari $t$. Prediksi return $\hat r_{t+1}$ dapat diubah menjadi
prediksi harga dengan
\begin{equation}
  \hat P_{t+1} = P_t \, e^{\hat r_{t+1}}.
  \label{eq:harga}
\end{equation}
Persamaan~\eqref{eq:harga} memperlihatkan mengapa grafik prediksi harga hampir selalu terlihat
akurat: $\hat P_{t+1}$ selalu dekat dengan $P_t$ karena $\hat r_{t+1}$ biasanya kecil. Oleh
karena itu, akurasi prediksi harga harus dibandingkan dengan prediksi naif
$\hat P_{t+1} = P_t$.

\subsection{Fitur}
Semua fitur pada hari $t$ hanya memakai informasi yang tersedia saat penutupan hari $t$.
Hal ini diverifikasi dengan pengujian otomatis: mengubah data setelah hari $t$ tidak boleh
mengubah nilai fitur sampai hari $t$.

\begin{table}[H]
  \centering
  \caption{Daftar fitur.}
  \label{tab:fitur}
  \small
  \begin{tabularx}{\linewidth}{lX}
    \toprule
    Fitur & Keterangan \\
    \midrule
    \texttt{ret\_lag1}--\texttt{ret\_lag5} & Log return hari ini sampai 4 hari sebelumnya \\
    \texttt{ret\_mean\_\{5,10,20\}} & Rata-rata return 5, 10, dan 20 hari terakhir (momentum) \\
    \texttt{vol\_\{5,10,20\}} & Simpangan baku return 5, 10, dan 20 hari terakhir (volatilitas) \\
    \texttt{vol\_ratio\_5\_20} & Rasio volatilitas 5 hari terhadap 20 hari \\
    \texttt{dist\_ma20}, \texttt{dist\_ma50} & Jarak harga dari rata-rata bergerak 20 dan 50 hari \\
    \texttt{rsi\_14} & \emph{Relative Strength Index} 14 hari \citep{wilder1978} \\
    \texttt{macd}, \texttt{macd\_hist} & MACD (EMA 12 $-$ EMA 26) dan histogramnya, dibagi harga \\
    \texttt{hl\_range} & Rentang harga harian $(H_t - L_t)/P_t$ \\
    \texttt{close\_in\_range} & Posisi harga penutupan dalam rentang harian \\
    \texttt{volume\_z20} & Skor-$z$ log volume terhadap 20 hari terakhir \\
    \texttt{day\_of\_week} & Hari dalam minggu (Senin = 0) \\
    \bottomrule
  \end{tabularx}
\end{table}

\subsection{Model}
Lima model dibandingkan. Tiga di antaranya adalah baseline yang tidak membutuhkan
\emph{machine learning}. Model yang baik harus mampu mengalahkan baseline tersebut.
\begin{enumerate}
  \item \textbf{Random walk}: $\hat r_{t+1} = 0$, artinya harga besok sama dengan harga hari ini.
  \item \textbf{Rata-rata historis}: $\hat r_{t+1}$ adalah rata-rata return pada data latih.
  \item \textbf{Momentum}: $\hat r_{t+1} = r_t$, artinya return hari ini terulang besok.
  \item \textbf{Regresi Ridge} \citep{hoerl1970}: regresi linear dengan penalti $L_2$
        ($\alpha = 50$) pada fitur yang sudah distandardisasi.
  \item \textbf{LightGBM} \citep{ke2017}: \emph{gradient boosting} berbasis pohon keputusan
        (300 pohon, \emph{learning rate} 0{,}02, maksimal 15 daun, minimal 50 sampel per daun,
        dengan \emph{subsampling} baris dan kolom 80\%). Parameter dipilih konservatif untuk
        mengurangi \emph{overfitting} pada data yang sangat berisik.
\end{enumerate}

\subsection{Validasi Walk-Forward}
Data deret waktu tidak boleh dibagi secara acak. Laporan ini memakai validasi
\emph{walk-forward} dengan jendela yang terus bertambah (\emph{expanding window}):
\begin{enumerate}
  \item Model pertama dilatih dengan seluruh data yang hari targetnya sebelum 1 Januari 2026
        (sekitar \JumlahLatih{} hari bursa).
  \item Model memprediksi 21 hari bursa berikutnya (sekitar satu bulan).
  \item Data 21 hari tersebut ditambahkan ke data latih, model dilatih ulang, lalu memprediksi
        21 hari berikutnya. Langkah ini diulang sampai akhir data.
\end{enumerate}
Dengan cara ini setiap prediksi pada tahun 2026 dibuat hanya dengan informasi yang
benar-benar tersedia pada saat itu, sama seperti jika model dipakai secara nyata.
Periode uji mencakup \JumlahHariUji{} hari bursa, dari \UjiAwal{} sampai \UjiAkhir.

\subsection{Metrik Evaluasi}
Untuk $n$ hari uji dengan return aktual $r_i$ dan prediksi $\hat r_i$:
\begin{align}
  \text{Akurasi arah} &= \frac{1}{n}\sum_{i=1}^{n} \mathbf{1}\!\left[\operatorname{sign}(\hat r_i) = \operatorname{sign}(r_i)\right], \\
  R^2_{\mathrm{RW}} &= 1 - \frac{\sum_i (r_i - \hat r_i)^2}{\sum_i r_i^2}, \\
  \text{MAE} &= \frac{1}{n}\sum_{i=1}^{n} |r_i - \hat r_i|, \qquad
  \text{RMSE}_{\text{harga}} = \sqrt{\frac{1}{n}\sum_{i=1}^{n} (P_i - \hat P_i)^2}.
\end{align}
Prediksi dianggap ``naik'' jika $\hat r_i > 0$. Hari ketika harga penutupan tidak berubah
($r_i = 0$) tidak dihitung dalam akurasi arah karena tidak memiliki arah. Hal ini penting
untuk saham di BEI, yang cukup sering ditutup di harga yang sama dengan hari sebelumnya.
Model random walk tidak pernah memberi arah, sehingga tidak memiliki nilai akurasi arah. $R^2_{\mathrm{RW}}$ adalah $R^2$
\emph{out-of-sample} terhadap random walk \citep{campbell2008}. Nilai positif berarti model
lebih baik daripada selalu menebak return 0\%, sedangkan nilai negatif berarti lebih buruk.
Untuk return saham harian, nilai $R^2_{\mathrm{RW}}$ di atas 1\% sudah tergolong sangat baik.

\subsection{Simulasi Strategi (Backtest)}
Prediksi diuji dengan strategi sederhana. Pada penutupan hari $t$, jika $\hat r_{t+1} > 0$
investor memegang saham sampai penutupan hari $t+1$. Jika tidak, dana disimpan tunai.
\emph{Short selling} tidak dipakai karena umumnya tidak tersedia bagi investor ritel di BEI.
Biaya transaksi sebesar 0{,}15\% untuk pembelian dan 0{,}25\% untuk penjualan (termasuk PPh
final 0{,}1\%) dikenakan setiap kali posisi berubah. Strategi dibandingkan dengan
\emph{buy \& hold}, yaitu membeli pada awal periode uji dan menahan saham sampai akhir.

Kinerja diukur dengan return total, rasio Sharpe tahunan
$\text{Sharpe} = \sqrt{252}\,\bar{x}/s_x$ dengan $x$ return harian strategi
\citep{sharpe1966}, dan \emph{drawdown} maksimum, yaitu penurunan terbesar dari puncak
nilai portofolio.

% ---------------------------------------------------------------------------
\section{Hasil}

\subsection{Ringkasan}
\RingkasanHasil

\begin{table}[H]
  \centering
  \caption{Ringkasan hasil periode uji 2026. Baseline terbaik adalah baseline dengan akurasi
  arah tertinggi untuk saham tersebut. RMSE naif memakai prediksi $\hat P_{t+1} = P_t$.}
  \label{tab:ringkasan}
  \small
  \input{generated/tabel_ringkasan.tex}
\end{table}

\newcommand{\bagiansaham}[2]{%
  \subsection{#2 (#1)}
  \input{generated/hasil_#1.tex}

  \begin{table}[H]
    \centering
    \caption{Hasil semua model untuk #1 pada periode uji. $^{\dagger}$Baseline.
    Akurasi arah, $R^2_{\mathrm{RW}}$, dan MAE mengukur prediksi return. Sharpe, return,
    \emph{drawdown}, dan jumlah transaksi mengukur strategi setelah biaya.}
    \small
    \resizebox{\linewidth}{!}{\input{generated/tabel_model_#1.tex}}
  \end{table}

  \begin{figure}[H]
    \centering
    \includegraphics[width=\linewidth]{prediksi_#1.pdf}
    \caption{#1 selama 2026. Atas: harga penutupan aktual dan prediksi LightGBM untuk satu
    hari ke depan. Bawah: return harian aktual (batang) dan prediksi LightGBM (garis).}
  \end{figure}
}

\bagiansaham{ASII}{Astra International}
\bagiansaham{BBRI}{Bank Rakyat Indonesia}
\bagiansaham{TLKM}{Telkom Indonesia}

\subsection{Kinerja Strategi}
Gambar~\ref{fig:ekuitas} memperlihatkan pertumbuhan Rp1 yang diinvestasikan pada awal 2026
untuk setiap strategi.

\begin{figure}[H]
  \centering
  \includegraphics[width=\linewidth]{ekuitas.pdf}
  \caption{Nilai Rp1 yang diinvestasikan pada awal periode uji, setelah biaya transaksi.}
  \label{fig:ekuitas}
\end{figure}

\subsection{Akurasi per Bulan}
Akurasi arah yang berada di sekitar 50\% secara keseluruhan bisa menyembunyikan periode
ketika model bekerja baik atau buruk. Gambar~\ref{fig:bulanan} memperlihatkan akurasi
LightGBM untuk setiap bulan di tahun 2026. Setiap bulan hanya berisi sekitar 20 hari bursa,
jadi akurasi bulanan antara 35\% dan 65\% masih wajar terjadi secara kebetulan.

\begin{figure}[H]
  \centering
  \includegraphics[width=\linewidth]{akurasi_bulanan.pdf}
  \caption{Akurasi arah LightGBM per bulan pada tahun 2026. Bulan dengan kurang dari 10 hari
  bursa pada periode uji tidak ditampilkan.}
  \label{fig:bulanan}
\end{figure}

\subsection{Fitur yang Paling Berpengaruh}
\begin{figure}[H]
  \centering
  \includegraphics[width=\linewidth]{fitur.pdf}
  \caption{Delapan fitur dengan kontribusi \emph{gain} terbesar pada model LightGBM yang
  dilatih dengan seluruh data.}
  \label{fig:fitur}
\end{figure}
Nilai \emph{importance} menunjukkan fitur yang paling sering dipakai model untuk memecah
data, tetapi tidak berarti fitur tersebut benar-benar memiliki daya prediksi. Pada data
yang sangat berisik, model juga dapat memakai fitur untuk menyesuaikan diri dengan
\emph{noise}.

\subsection{Prediksi Hari Bursa Berikutnya}
Model LightGBM yang dilatih dengan seluruh data menghasilkan prediksi untuk hari bursa
setelah data terakhir pada Tabel~\ref{tab:prediksi}. Mengingat hasil evaluasi di atas,
prediksi ini sebaiknya dibaca sebagai ilustrasi cara kerja model, bukan rekomendasi
investasi.

\begin{table}[H]
  \centering
  \caption{Prediksi LightGBM untuk hari bursa berikutnya. Sinyal \textsc{beli} jika
  prediksi return positif, \textsc{tunai} jika tidak.}
  \label{tab:prediksi}
  \small
  \input{generated/tabel_prediksi.tex}
\end{table}

% ---------------------------------------------------------------------------
\section{Pembahasan}

\paragraph{Prediksi harga terlihat akurat, tetapi itu menyesatkan.}
Pada grafik harga, garis prediksi hampir menempel pada garis aktual. Hal ini bukan bukti
bahwa model pintar, melainkan akibat Persamaan~\eqref{eq:harga}: prediksi selalu dimulai dari
harga hari ini. Perbandingan yang jujur adalah dengan prediksi naif, dan pada
Tabel~\ref{tab:ringkasan} terlihat bahwa selisih RMSE keduanya sangat kecil.

\paragraph{Return harian didominasi \emph{noise}.}
Prediksi return dari model hampir selalu jauh lebih kecil daripada return aktual (panel bawah
pada grafik prediksi). Model yang dilatih dengan galat kuadrat akan menyusutkan prediksinya
ke arah rata-rata ketika sinyal lemah. Hal ini konsisten dengan hipotesis pasar efisien dan
dengan literatur yang menemukan bahwa prediksi return \emph{out-of-sample} sangat sulit
\citep{welch2008}.

\paragraph{Biaya transaksi menggerus keunggulan kecil.}
Setiap kali strategi berpindah antara saham dan tunai, biaya sekitar 0{,}4\% untuk satu kali
beli dan jual harus ditutup oleh keuntungan. Akurasi arah yang hanya sedikit di atas 50\%
tidak cukup untuk menutup biaya tersebut jika strategi sering berganti posisi.

\paragraph{Keterbatasan.}
\begin{itemize}
  \item Periode uji tidak lebih dari satu tahun (\JumlahHariUji{} hari bursa per saham),
        sehingga perbedaan akurasi beberapa poin persen masih bisa terjadi secara kebetulan.
        Sebagai gambaran, simpangan baku akurasi tebakan acak untuk \JumlahHariUji{} hari
        adalah sekitar \SimpanganAcak{} poin persen.
  \item Simulasi mengasumsikan transaksi terjadi tepat di harga penutupan tanpa
        \emph{slippage}, serta mengabaikan batas fraksi harga dan ukuran lot.
  \item Hiperparameter model tidak dioptimalkan. Hal ini disengaja untuk menghindari
        \emph{overfitting} pada periode uji, tetapi berarti masih ada ruang untuk perbaikan.
  \item Fitur hanya berasal dari harga dan volume. Informasi lain seperti laporan keuangan,
        berita, suku bunga, atau arus dana asing tidak digunakan.
\end{itemize}

% ---------------------------------------------------------------------------
\section{Kesimpulan dan Saran}

\subsection{Kesimpulan}
\input{generated/kesimpulan.tex}

Secara umum, hasil ini mendukung pandangan bahwa return harian saham besar di BEI sulit
diprediksi hanya dengan data harga dan volume. Evaluasi yang jujur dengan validasi
\emph{walk-forward}, perbandingan dengan baseline, dan perhitungan biaya transaksi sangat
penting. Tanpa ketiganya, model yang tidak berguna dapat terlihat sangat akurat.

\subsection{Saran}
\begin{enumerate}
  \item Menambahkan fitur non-harga seperti sentimen berita (misalnya dengan IndoBERT),
        data makro, atau arus dana asing.
  \item Memprediksi volatilitas alih-alih arah harga, karena volatilitas umumnya lebih mudah
        diprediksi dan berguna untuk manajemen risiko.
  \item Memperpanjang periode uji, misalnya 2022--2026, agar kesimpulan lebih kuat secara
        statistik.
  \item Membandingkan dengan model \emph{deep learning} seperti LSTM atau Transformer
        dengan prosedur validasi yang sama.
\end{enumerate}

\vspace{1em}
\noindent\small\textit{Penafian: laporan ini dibuat untuk tujuan pembelajaran dan portofolio,
bukan merupakan saran atau rekomendasi investasi.}
\normalsize

% ---------------------------------------------------------------------------
\begin{thebibliography}{9}
\bibitem[Campbell \& Thompson(2008)]{campbell2008}
  Campbell, J.~Y., \& Thompson, S.~B. (2008). Predicting excess stock returns out of sample:
  Can anything beat the historical average? \emph{The Review of Financial Studies}, 21(4),
  1509--1531.
\bibitem[Fama(1970)]{fama1970}
  Fama, E.~F. (1970). Efficient capital markets: A review of theory and empirical work.
  \emph{The Journal of Finance}, 25(2), 383--417.
\bibitem[Hoerl \& Kennard(1970)]{hoerl1970}
  Hoerl, A.~E., \& Kennard, R.~W. (1970). Ridge regression: Biased estimation for
  nonorthogonal problems. \emph{Technometrics}, 12(1), 55--67.
\bibitem[Ke et al.(2017)]{ke2017}
  Ke, G., Meng, Q., Finley, T., Wang, T., Chen, W., Ma, W., Ye, Q., \& Liu, T.-Y. (2017).
  LightGBM: A highly efficient gradient boosting decision tree. \emph{Advances in Neural
  Information Processing Systems}, 30.
\bibitem[Sharpe(1966)]{sharpe1966}
  Sharpe, W.~F. (1966). Mutual fund performance. \emph{The Journal of Business}, 39(1),
  119--138.
\bibitem[Welch \& Goyal(2008)]{welch2008}
  Welch, I., \& Goyal, A. (2008). A comprehensive look at the empirical performance of
  equity premium prediction. \emph{The Review of Financial Studies}, 21(4), 1455--1508.
\bibitem[Wilder(1978)]{wilder1978}
  Wilder, J.~W. (1978). \emph{New Concepts in Technical Trading Systems}. Trend Research.
\end{thebibliography}

% ---------------------------------------------------------------------------
\appendix
\section{Kode Program}
\label{sec:kode}

\subsection{Cara Menjalankan}
\begin{lstlisting}[language=bash,numbers=none]
# 1. Instal paket
pip install -e ".[download]"

# 2. Unduh data ke Excel (menghasilkan data/data_saham.xlsx)
python scripts/download_data.py

# 3. Hitung hasil: tabel, grafik, dan paragraf (ke laporan/generated/)
python -m stockpred.report --excel data/data_saham.xlsx

# 4. Kompilasi laporan PDF
cd laporan && latexmk -pdf laporan.tex
\end{lstlisting}

\subsection{Unduh Data ke Excel (\texttt{scripts/download\_data.py})}
\lstinputlisting{../scripts/download_data.py}

\subsection{Pengaturan (\texttt{config.py})}
\lstinputlisting{../src/stockpred/config.py}

\subsection{Membaca Data (\texttt{data.py})}
\lstinputlisting{../src/stockpred/data.py}

\subsection{Fitur (\texttt{features.py})}
\lstinputlisting{../src/stockpred/features.py}

\subsection{Model (\texttt{models.py})}
\lstinputlisting{../src/stockpred/models.py}

\subsection{Validasi Walk-Forward dan Metrik (\texttt{evaluate.py})}
\lstinputlisting{../src/stockpred/evaluate.py}

\subsection{Backtest (\texttt{backtest.py})}
\lstinputlisting{../src/stockpred/backtest.py}

\subsection{Pembuatan Tabel dan Grafik (\texttt{report.py})}
\lstinputlisting{../src/stockpred/report.py}

\end{document}
